\section{Conclusão}
\label{sec-conclusao}

Neste trabalho, aplicou-se a meta-heurística \textit{Ant Colony System} (ACS) para resolver o problema de acasalamento otimizado em rebanhos, equilibrando simultaneamente dois objetivos conflitantes: minimizar a coancestralidade entre os descendentes e maximizar o ganho genético esperado. A formulação adotada modelou o problema como um grafo bipartido, no qual cada formiga constrói iterativamente uma solução orientada tanto pelo nível de feromônio quanto pela heurística genética, sempre respeitando as restrições reprodutivas impostas aos machos. 

Os experimentos realizados, baseados na média de 30 execuções independentes por cenário, evidenciaram que o aumento do limite de uso dos machos conduz a um ganho genético crescente, porém acompanhado de um aumento igualmente progressivo da coancestralidade. Assim, nenhum limite testado se mostrou simultaneamente ótimo para ambos os critérios. A normalização dos resultados e a aplicação de uma função de utilidade permitiram quantificar o equilíbrio entre custo (coancestralidade) e benefício (mérito genético), indicando que o limite de 15 acasalamentos por macho fornece o melhor compromisso global, oferecendo ganho genético expressivo sem comprometer excessivamente a diversidade.

Conclui-se, portanto, que o ACS é uma abordagem adequada para problemas multicritério de otimização genética no contexto de manejo animal. Além disso, o estudo reforça que a definição do limite de uso dos reprodutores deve considerar não apenas o mérito genético imediato, mas também a sustentabilidade da variabilidade genética ao longo das gerações. Como trabalhos futuros, sugere-se a utilização de funções de utilidade adaptativas, a incorporação de parâmetros econômicos reais, a extensão para múltiplas gerações e a comparação direta com outras meta-heurísticas, como Algoritmos Genéticos e Simulated Anneling, a fim de aprofundar o potencial do ACS em programas de melhoramento animal.
